\chapter{Accuracy: The Independent Evidence}
\label{ch:accuracy}

\section{Peer-reviewed validation}

\textbf{Leach, Forrester, Mears \& Roberts (2017)}~\cite{leach2017}
remains the key criterion study: 240 shots (driver, 7-iron, wedge)
measured simultaneously by a TrackMan Pro~IIIe, a Foresight GC2+HMT, and
a four-camera 5,400\,fps optical reference. Findings: ball parameters
agreed well on both devices (TrackMan clubhead-speed median difference
$-0.4$\,mph; ball speed $+0.2$\,mph; launch angle $0.0\degs$), but
\textbf{club parameters were materially worse for both} --- the authors
endorse ball data for research use and advise caution on club data. This
is the empirical face of the directness hierarchy of
\cref{sec:hierarchy}: what is measured agrees; what is derived diverges.

\textbf{TrackMan 4 indoor reliability (J. Sports Sciences,
2024)}~\cite{tm4reliability}: within- and between-session reliability in
high-level golfers indoors --- ICC 0.99 for club speed, 0.97--0.99 for
ball speed, but \textbf{spin-rate ICC as low as 0.02--0.60}: even the
reference radar's indoor spin channel is fragile when flight is
truncated.

A 2025 Mevo+ vs.\ TrackMan~4 indoor agreement study
exists~\cite{mevoplusstudy} (abstract-level access only at review time),
extending the same pattern down-market.

\section{Robot and industry testing}

Golf Laboratories robot testing (reported via Foresight and third
parties) puts GCQuad center-strike spin standard deviation at
$\approx$82\,rpm versus $\approx$175\,rpm for TrackMan~4 --- the
photometric spin advantage in its clearest form~\cite{golflabsrobot}.
MyGolfSpy's indoor comparison against a GCQuad reference found the
camera-assisted MLM2PRO among the tightest budget units for launch,
spin, and carry, with radar-only units (R10, original Mevo) trailing
indoors~\cite{mygolfspyindoor}. Manufacturer-quoted specs bracket the
market: $\pm1$\,mph ball speed and $\pm1\degs$ launch angles are common
claims; club-face claims ($\pm2\degs$, ``calculated'') are visibly
weaker.

\section{What drives inter-device disagreement}

Synthesizing the validation literature and the architecture analysis:

\begin{enumerate}
\item \textbf{Reference-point differences} in club speed (geometric
center vs.\ fastest return) --- several mph of systematic spread.
\item \textbf{Derived face data}: D-plane inversion amplifies
launch-direction bias by $\sim$1.15$\times$ and cannot see gear effect,
so radar face angle diverges from optical face angle most on off-center
strikes.
\item \textbf{Indoor spin}: sideband SNR collapses on short flights and
clean balls; systems silently switch to estimation, and estimated spin
feeds the carry model.
\item \textbf{Flight-model differences}: identical launch conditions
produce different simulated carries across vendors; this is a modeling
disagreement, not a measurement one.
\item \textbf{Alignment}: a unit misaligned to the target line biases
path, face, and launch direction coherently --- the cheapest error to fix
and the most common in practice.
\end{enumerate}

\begin{implication}
For OpenFlight validation against the MLM2PRO (the project's reference
instrument): compare ball speed and launch angles directly; compare spin
only on RPT/RCT marked balls where the MLM2PRO's measurement is
camera-based; expect club-speed offsets from reference-point differences
and calibrate a per-club correction rather than chasing agreement; and
log raw I/Q captures so spin-detection improvements can be replayed
against historical shots.
\end{implication}
